% !TEX encoding = UTF-8 Unicode

% ------------------------------------------------------------------------------------------------------
%	Formatvorlage für wissenschaftliche Arbeiten (Diplomarbeit, Bachelorarbeit, Masterarbeit, Seminararbeit, Projektarbeit, Dokumentation)
% ------------------------------------------------------------------------------------------------------
%	ursprünglich erstellt von Stefan Macke, 24.04.2009
%	http://blog.stefan-macke.de
%
%	erweitert von Felix Rupp, 21.05.2013
%	http://www.felixrupp.com/
%
%
%
% erweitert und angepasst von Moritz Rupp
% moritz.rupp@gmail.com
%
% Version 1.0.2
% Datum: 12.07.2013


% Dokumentenkopf ---------------------------------------------------------------------------------------
%   Diese Vorlage basiert auf "scrreprt" aus dem koma-script.
% ------------------------------------------------------------------------------------------------------
\documentclass[
    12pt, % Schriftgröße
    DIV10, % Änderung der Größe des Satzspiegels (bedruckbarer Bereich einer Seite), nur in Verbindung mit koma-script verwendbar
    ngerman, % für Umlaute, Silbentrennung etc.
    a4paper, % Papierformat
    oneside, % einseitiges Dokument
    titlepage, % es wird eine Titelseite verwendet
    parskip=half, % Abstand zwischen Absätzen (halbe Zeile)
    headings=normal, % Größe der Überschriften verkleinern
    listof=totoc, % Verzeichnisse im Inhaltsverzeichnis aufführen
    bibliography=totoc, % Literaturverzeichnis im Inhaltsverzeichnis aufführen
    index=totoc, % Index im Inhaltsverzeichnis aufführen
    captions=tableheading, % Beschriftung von Tabellen unterhalb ausgeben
    final % Status des Dokuments (final/draft)
]{scrreprt}

% UTF8 und T1 Fontencoding -----------------------------------------------------------------------------
\usepackage[utf8]{inputenc}
\usepackage[T1]{fontenc}


% Meta-Informationen -----------------------------------------------------------------------------------
%   Informationen über das Dokument, wie z.B. Titel, Autor, Matrikelnr. etc
%   werden in der Datei Meta.tex definiert und können danach global
%   verwendet werden.
% ------------------------------------------------------------------------------------------------------
% !TEX encoding = UTF-8 Unicode
% !TEX root =  Paper.tex

% Meta information ------------------------------------------------------------------------------------
% Definiton of global parameters, they can be used in the whole document (e.g. on the titlepage)
% -------------------------------------------------------------------------------------------------------
\newcommand{\myTitle}{The Effects of National Business Systems on HRM in Multinational Companies} % the title for text
\newcommand{\myTitleTitlepage}{The Effects of National Business Systems on HRM in \\ Multinational Companies} % the title for the titlepage

\newcommand{\mySubtitle}{} 
\newcommand{\mySubtitleTitlepage}{}

\newcommand{\study}{General Management}
\newcommand{\area}{Faculty of Economic and Social Sciences}

\newcommand{\myAuthor}{Michaela Wunder}

\newcommand{\keywords}{}

\newcommand{\born}{Ellwangen}

\newcommand{\firstreviewer}{Professor Dr.\,Kerstin Pull}
\newcommand{\secondreviewer}{}

\newcommand{\submissionDate}{09/04/2015}

\newcommand{\university}{Eberhard Karls Universität}
\newcommand{\place}{Tübingen}

\newcommand{\creator}{texmaker 4.2}

\newcommand{\bibliographyheading}{Reference list}

\newcommand{\TableCellPadding}{1.2}


% benötigte Packages -----------------------------------------------------------------------------------
%   LaTeX-Packages, die benötigt werden, sind in die Datei Packages.tex
%   "ausgelagert", um diese Vorlage möglichst übersichtlich zu halten.
% ------------------------------------------------------------------------------------------------------
% !TEX encoding = UTF-8 Unicode
% !TEX root =  Paper.tex

% Customisation of the page layout ---------------------------------------------------
%   see Layout.tex
% -----------------------------------------------------------------------------------------
\usepackage[
    automark, % automatically create chapter name in the headline
    headsepline, % seperation line beneath headline
    ilines % left seperate line
]{scrpage2}

% Hacks in order to prevent warnings
\usepackage{scrhack}

% Customisation of the language -------------------------------------------------
\usepackage[british]{babel}


% Special character 
% -----------------------------------------------------------------------------------------
\usepackage{textcomp} % Euro-sign etc.


% Font --------------------------------------------------------------------------------
\usepackage{lmodern} % better fonts
\usepackage{relsize} % font size relatively


% Better underlining ---------------------------------------------
\usepackage[normalem]{ulem}


% Graphics -----------------------------------------------------------------------------
% Enable the inclusion of JPEG
\usepackage[dvips,final]{graphicx}
% Path to the graphics
\graphicspath{{Images/}}
\usepackage{picture}


% Commands from AMSTeX for mathematical symbols, e.g.  \boldsymbol \mathbb
\usepackage{amsmath,amsfonts}


% Use Euro sign ------------------------------------------------------------
\usepackage{eurosym}


% For index with \printindex -----------------------------------------------
\usepackage{makeidx}


% Easy definition of margins and line spacing etc. ------------
\usepackage{setspace}
\usepackage{geometry}

% Floating figures ---------------------------------------------------------
\usepackage[vflt]{floatflt}
\usepackage{subfig}
\usepackage{tocbasic}

% Inclusion of program code -----------------------------------------------
\usepackage{listings}
\usepackage{xcolor} 
\definecolor{hellgelb}{rgb}{1,1,0.9}
\definecolor{colKeys}{rgb}{0.8,0,0.5}
\definecolor{colIdentifier}{rgb}{0.6,0,0.3}
\definecolor{colComments}{rgb}{0,0.5,0}
\definecolor{colString}{rgb}{0,0,1}

\lstset{
    float=htbp,
    basicstyle=\ttfamily\color{black}\small\smaller,
    identifierstyle=,%\color{colIdentifier},
    keywordstyle=\color{colKeys}\bfseries,
    stringstyle=\color{colString},
    commentstyle=\color{colComments},
    columns=flexible,
    tabsize=4,
    frame=single,
    extendedchars=true,
    showspaces=false,
    showstringspaces=false,
    numbers=left,
    numberstyle=\tiny,
    breaklines=true,
    backgroundcolor=\color{hellgelb},
    breakautoindent=true,
    escapeinside={(*}{*)},
    literate={Ö}{{\"O}}1 {Ä}{{\"A}}1 {Ü}{{\"U}}1 {ß}{{\ss}}2 {ü}{{\"u}}1 {ä}{{\"a}}1 {ö}{{\"o}}1 {µ}{\textmu}1
 }

% --------------------------------------------------------------------------------------------
% 
% Definition of source code
%
% Sample Defininition of CSS:
\lstdefinelanguage{CSS}{
    morestring=[b]',
    morestring=[b]",
    comment=[l]{/*}{*/},
    sensitive=false,
    morekeywords={accelerator,adjust,after,align,attachment,azimuth,background,before,behavior,binding,border,bottom,bottomright,bottomleft,break,caption-side,char,clear,clip,color,colors,cue,cursor,collapse,decoration,direction,display,elevation,empty-cells,family,filter,float,flow,focus,font,grid,height,image,increment,indent,input,inside,ime-mode,include-source,justify,last,layer,left,letter,line,list,margin,marker,marks,max,min,mode,modify,moz,offsett,opacity,orphans,outline,overflow,overflow-X,overflow-Y,overhang,position,position-x,position-y,padding,page,pause,pitch,play-during,position,quotes,radius,range,repeat,replace,reset,richness,right,ruby,set-link-source,select,shadow,size,spacing,speak,speak-header,speak-numeral,speak-punctuation,speech-rate,stress,stretch,style,table-layout,text,transform,text-autospace,text-kashida-space,top,topleft,topright,type,underline,unicode-bidi,use-link-source,user,variant,vertical,visibility,voice-family,volume,white-space,weight,widows,width,word,wrap,word-wrap,writing-mode,z-index,zoom}
}

% link URL, break long URL etc. -------------------------------------
\usepackage{url}


% include biblatex -----------------------------------------------------------------------
\usepackage[
	citestyle=my-authoryear-ibid,
	bibstyle=authoryear,
	backend=biber,
	babel=hyphen,
	urldate=comp,
	dateabbrev=false,
	maxcitenames=1
]{biblatex}
\usepackage[english=british]{csquotes}


% PDF Options -------------------------------------------------------------------------
\usepackage[
    bookmarks,
    bookmarksopen=true,
    colorlinks=true,
% use this colour definitions for links with colour
    linkcolor=red, % simple internal link
    anchorcolor=black,% anchor text
    citecolor=blue, % Anchor to bibliography entries in the text
    filecolor=magenta, % Anchors to local files
    menucolor=red, % Acrobat menu entries
    urlcolor=cyan, 
%
% use this colour definitions for pring: everything is black
%
    %linkcolor=black, % simple internal link
    %anchorcolor=black, % anchor text
    %citecolor=black, %Anchor to bibliography entries in the text
    %filecolor=black, % Anchors to local files
    %menucolor=black, % Acrobat menu entries
    %urlcolor=black, 
%
    plainpages=false, % for the correct creation of bookmarks
    pdfpagelabels, % for the correct creation of bookmarks
    hypertexnames=false, % for the correct creation of bookmarks
    linktocpage % link the page numbers instead of the whole toc entry
]{hyperref}

% Commands with print special characters cause errors if they are used with hyperref
\hypersetup{
    pdftitle={\myTitle \mySubtitle},
    pdfauthor={\myAuthor},
    pdfcreator={\creator},
    pdfsubject={\myTitle \mySubtitle},
    pdfkeywords={\keywords},
}

% For the glossary and abbreviations -------------------------
\usepackage[
	%xindy,
	nonumberlist,	% keine Seitenzahlen anzeigen
	acronym,			% ein Abkürzungsverzeichnis erstellen
	toc,					% Einträge im Inhaltsverzeichnis
	shortcuts			% acronym shortcuts
]{glossaries}

% Package for a clean inclusion of external PDF files -----------------
\usepackage[final]{pdfpages}

% Continuous numeration of footnotes -------------------------------
\usepackage{chngcntr}


% nicer tables --------------------------------------------------------------------
\usepackage{tabularx}
\usepackage{multirow}


% for long tables ---------------------------------------------------------------------
\usepackage{longtable}
\usepackage{ltxtable}
\usepackage{filecontents}
\usepackage{array}
\usepackage{ragged2e}
\usepackage{lscape}


% rotation of elements -------------------------------------------------------
\usepackage{rotating}


% Change the formatting of lists --------------------------------------------------
\usepackage{paralist}
\usepackage{enumitem}


% Necessary for the creation of own commands -------------------------------------
\usepackage{ifthen}
\usepackage{forloop}


% defines commands for \todo and \listoftodos --------------------------------
\usepackage[english]{todonotes}

% Enables strikethrough for text
\usepackage{ulem}


% used for the end of macros without parameter
\usepackage{xspace}

% used for bold labels in bibliography entries
\usepackage{xpatch}

% for multiple columns and rows
\usepackage{multicol}
\usepackage{multirow}

% better hyphenation etc.
\usepackage{microtype}

% This is used for the dummy text
\usepackage{blindtext}

%% Uncomment this, if you want to let chapters start at the current position and not on a new page. ------------------------------------------------------
%\usepackage{etoolbox}
%\makeatletter
%\patchcmd{\chapter}{\if@openright\cleardoublepage\else\clearpage\fi}{}{}{}
%\makeatother


% Erstellung eines Index und Abkürzungsverzeichnisses/Glossars aktivieren ------------------------------
\makeindex{}
%\makenomenclature{}


% Kopf- und Fußzeilen, Seitenränder etc. ---------------------------------------------------------------
% !TEX encoding = UTF-8 Unicode
% !TEX root =  Arbeit.tex

% Zeilenabstand 1,5 Zeilen ---------------------------------------------------------------------
\onehalfspacing{}


% Seitenränder ------------------------------------------------------------------------------------
\setlength{\topskip}{\ht\strutbox} % behebt Warnung von geometry
\geometry{paper=a4paper,left=35mm,right=35mm,top=35mm}


% Kopf- und Fußzeilen --------------------------------------------------------------------------
\pagestyle{scrheadings}

% Kopf- und Fußzeile auch auf Kapitelanfangsseiten
\renewcommand*{\chapterpagestyle}{scrheadings} 

% Schriftform der Kopfzeile
\renewcommand{\headfont}{\normalfont}


% Kopfzeile für einseitiges Dokument:
\ihead{%\small{\textsc{\titel}	\\
	%\untertitel
	%\\[2ex]
	\textit{\headmark}%}
}
\chead{}
\rohead{\includegraphics[scale=0.06]{\logo}}
\setlength{\headheight}{21mm} % Höhe der Kopfzeile
% Kopfzeile über den Text hinaus verbreitern
\setheadwidth[0pt]{textwithmarginpar} 
\setheadsepline[text]{0.4pt} % Trennlinie unter Kopfzeile


%% Kopfzeile für Zweiseitiges Dokument:
%\ihead{\textit{\leftmark}}
%\chead{}
%\ohead{}
%\lehead{\includegraphics[scale=0.25]{\logo}}
%\rohead{\includegraphics[scale=0.25]{\logo}}
%\setheadsepline[text]{0.4pt} % Trennlinie unter Kopfzeile


% Fußzeile
%\ifoot{\copyright\ \autor}
\cfoot{}
\ofoot{\pagemark \\[4ex]}


% sonstige typographische Einstellungen ---------------------------------------------------
% erzeugt ein wenig mehr Platz hinter einem Punkt
\frenchspacing{}

% Schusterjungen und Hurenkinder vermeiden
\clubpenalty = 10000
\widowpenalty = 10000 
\displaywidowpenalty = 10000

% Quellcode-Ausgabe formatieren
\lstset{numbers=left, numberstyle=\tiny, numbersep=5pt, breaklines=true}
\lstset{emph={square}, emphstyle=\color{red}, emph={[2]root,base}, emphstyle={[2]\color{blue}}}

% Fußnoten fortlaufend durchnummerieren
\counterwithout{footnote}{chapter}




% eigene Definitionen für Silbentrennung ---------------------------------------------------------------
% !TEX encoding = UTF-8 Unicode
% !TEX root =  Arbeit.tex

% Trennvorschläge im Text werden mit \" angegeben
% untrennbare Wörter und Ausnahmen von der normalen Trennung können in dieser
% Datei mittels \hyphenation definiert werden

% \hyphenation{Mit-ar-bei-ter}


% eigene LaTeX-Befehle ---------------------------------------------------------------------------------
% !TEX encoding = UTF-8 Unicode
% !TEX root =  Arbeit.tex
% Eigene Befehle und typographische Auszeichnungen für diese


% einfaches Wechseln der Schrift, z.B.: \changefont{cmss}{sbc}{n} ---------------------------------------
\newcommand{\changefont}[3]{\fontfamily{#1} \fontseries{#2} \fontshape{#3} \selectfont}


% Abkürzungen mit korrektem Leerraum --------------------------------------------------------------------
\newcommand{\ua}{\mbox{u.\,a.\ }}
\newcommand{\zB}{\mbox{z.\,B.\ }}
\newcommand{\dahe}{\mbox{d.\,h.\ }}
\newcommand{\Vgl}{Vgl.\ }
\newcommand{\bzw}{bzw.\ }
\newcommand{\evtl}{evtl.\ }

\newcommand{\Abbildung}[1]{Abbildung~\ref{fig:#1}}

\newcommand{\bs}{$\backslash$}


% erzeugt ein Listenelement mit fetter Überschrift ------------------------------------------------------
\newcommand{\itemd}[2]{\item{\textbf{#1}}\\{#2}}


% einige Befehle zum Zitieren ---------------------------------------------------------------------------
\newcommand{\Zitat}[2][\empty]{\ifthenelse{\equal{#1}{\empty}}{\citep*{#2}}{\citep*[#1]{#2}}}
\newcommand{\ZitatSpezialA}[3]{\citep*[#1][#2]{#3}}
\newcommand{\ZitatSpezialB}[2]{\citep*[#1][]{#2}}


% zum Ausgeben von Autoren
\newcommand{\AutorName}[1]{\textsc{#1}}
\newcommand{\Autor}[1]{\AutorName{\citeauthor{#1}}}


% verschiedene Befehle um Wörter semantisch auszuzeichnen -----------------------------------------------
\newcommand{\NeuerBegriff}[1]{\textbf{#1}}

%\newcommand{\Fachbegriff}[2][\empty]{\ifthenelse{\equal{#1}{\empty}}{\textit{#2}\xspace}{\textit{#2}\xspace\footnote{#1}\nomenclature{#2}{#1}}}

%\newcommand{\FachbegriffSpezialA}[4]{\textit{#4}\footnote{#3}\label{fn:#1}\nomenclature{#4}{#2. Siehe auch Fu{\ss}zeile auf Seite~\pageref{fn:#1}.}}

%\newcommand{\FachbegriffSpezialB}[5]{\textit{#5}\footnote{#3}\label{fn:#1}\nomenclature{#4}{#2. Siehe auch Fu{\ss}zeile auf Seite~\pageref{fn:#1}.}}


% Beträge mit Währung -----------------------------------------------------------------------------------
\newcommand{\Betrag}[2][general]{#2\,\ifthenelse{\equal{#1}{dollar}}{\$}{}\ifthenelse{\equal{#1}{euro}}{€}{}\ifthenelse{\equal{#1}{yen}}{¥}{}\ifthenelse{\equal{#1}{cent}}{¢}{}\ifthenelse{\equal{#1}{pound}}{£}{}\ifthenelse{\equal{#1}{peso}}{₱}{}\ifthenelse{\equal{#1}{baht}}{฿}{}\ifthenelse{\equal{#1}{franc}}{₣}{}\ifthenelse{\equal{#1}{lira}}{₤}{}\ifthenelse{\equal{#1}{drachma}}{₯}{}\ifthenelse{\equal{#1}{pfennig}}{₰}{}\ifthenelse{\equal{#1}{general}}{¤}{}}


% Sonstiges ---------------------------------------------------------------------------------------------
\newcommand{\Eingabe}[1]{\texttt{#1}}
\newcommand{\Code}[1]{\texttt{#1}}
\newcommand{\Datei}[1]{\texttt{#1}}

\newcommand{\Datentyp}[1]{\textsf{#1}}
\newcommand{\XMLElement}[1]{\textsf{#1}}
\newcommand{\Webservice}[1]{\textsf{#1}}


% Beschriftung von Tabellen und Bildern ändern ----------------------------------------------------------
\addto\captionsngerman{
	\renewcommand{\figurename}{Abb.}
	\renewcommand{\tablename}{Tab.}
}


% Spaltendefinition rechtsbündig mit definierter Breite -------------------------------------------------
\newcolumntype{w}[1]{>{\raggedleft\hspace{0pt}}p{#1}}


% Linksbündige Tabellenspalten mit tabularx -------------------------------------------------------------
\newcolumntype{y}[1]{>{\RaggedRight\arraybackslash\hsize=#1\hsize}X}



% Literaturverzeichnis ---------------------------------------------------------------------------------
%   Das Literaturverzeichnis wird aus der BibTeX-Datenbank "Bibliographie.bib"
%   erstellt.
% ------------------------------------------------------------------------------------------------------
\bibliography{Bibliographie} % Aufruf: biblatex Bibliographie
\defbibheading{literaturverzeichnis}{\chapter*{\bibliographyheading}}


% Glossar --------------------------------------------------------------------------------
\makeglossaries
\loadglsentries{Inhalt/Glossar.tex}

% Hinzufügen der nicht benutzten Glossareinträge (falls noch eine Abkürzung vorhanden ist und nur diese verwendet wird)
%\glsadd{beispiel}


% Das eigentliche Dokument -----------------------------------------------------------------------------
%   Der eigentliche Inhalt des Dokuments beginnt hier. Die einzelnen Seiten
%   und Kapitel werden in eigene Dateien ausgelagert und hier nur inkludiert.
% ------------------------------------------------------------------------------------------------------
\begin{document}


% auch subsubsections nummerieren ----------------------------------------------------------------------
\setcounter{secnumdepth}{3}
% Nummerierungsebenen im Inhaltsverzeichnis
\setcounter{tocdepth}{2}


% Deckblatt und Abstract ohne Seitenzahl ---------------------------------------------------------------
\ofoot{}
% !TEX encoding = UTF-8 Unicode
% !TEX root =  Arbeit.tex

\thispagestyle{plain}
\begin{titlepage}

\changefont{cmss}{bx}{n}

\begin{center}

\LARGE{\textbf{\titel}}\\[1.5ex]
\Large{\textbf{\untertitelDeckblatt}}\\[6ex]
\Large{\textbf{\art}}\\[1.5ex]

\changefont{cmss}{m}{n}
\large{\fachgebiet \studienbereich}\\[6ex]

von

\large{\textbf{\autor}} \\[3ex]

September \jahr \\[6ex]

% 2 Logos -------------------
%\begin{tabular}{p{40mm}cp{40mm}c}
%	\includegraphics[scale=0.16]{Logo.pdf} & \includegraphics[trim=0mm 100mm 0mm 100mm, scale=0.25]{Logo.pdf}
%\end{tabular}

% 1 Logo ----------------------
\includegraphics[width=0.7\textwidth]{Logo.pdf}
\\[6ex]

\normalsize{}
\begin{tabular}{ll}
Bearbeitungszeitrum:  & \quad 12 Wochen\\
Matrikelnummer, Kurs: & \quad \matrikelnr, \kurs\\
Ausbildungsfirma: & \quad \firma, Firmenstandort\\
Betreuer:  & \quad \erstgutachter\\
Gutachter: & \quad \zweitgutachter\\
eingereicht in: & \quad \ort, am 12.\,August\,2011
\end{tabular}

\end{center}


% Die folgenden Kommentare entfernen um einen Sperrvermerk auf der 2. Seite zu erzeugen
%\singlespacing

%\newpage

%\thispagestyle{empty}

%\changefont{cmr}{m}{n}

%\small
%\noindent Die nachfolgende \art\ \textbf{enth\"alt vertrauliche Daten der Musterfabrik GmbH \& Co. Betriebs KG}. Ver\"offentlichungen oder Vervielf\"altigungen -- auch auszugsweise -- sind ohne ausdr\"uckliche Genehmigung der Musterfabrik GmbH \& Co. Betriebs KG nicht gestattet. Die \art\ ist nur den Erst- und Zweitgutachtern und dem Prüfungsausschuss zug\"anglich zu machen.\\[5ex]

\end{titlepage}


% Sperrvermerk ----------------------------------------------------------------------------
% !TEX encoding = UTF-8 Unicode
% !TEX root =  Paper.tex

\chapter*{Sperrvermerk}

\thispagestyle{empty}


Die vorliegende {\art} mit dem Titel \textit{\titel} beinhaltet interne und vertrauliche Informationen der \firma{}, sie ist daher mit einem Sperrvermerk versehen und wird ausschließlich zu Pr\"ufungszwecken im Studiengang {\studienbereich} der \hochschule{} \ort{} vorgelegt.

Jede Einsichtnahme und Ver\"offentlichung -- auch von Teilen der Arbeit oder in digitaler Form –- bedarf der vorherigen schriftlichen Zustimmung durch die {\firma}.

% Selbständigkeitserklärung ----------------------------------------------------------------------------
% !TEX encoding = UTF-8 Unicode
% !TEX root =  Paper.tex

\chapter*{Eidesstattliche Versicherung}

\thispagestyle{empty}

Ich, \autor, Matrikel-Nr.\ \matrikelnr, versichere hiermit, dass ich meine \art\xspace mit dem Thema
\begin{quote}
\textit{\titel} \textit{\untertitel}
\end{quote}
selbständig verfasst und keine anderen als die angegebenen Quellen und Hilfsmittel benutzt habe, wobei ich alle wörtlichen und sinngemäßen Zitate als solche gekennzeichnet habe. Die Arbeit wurde bisher keiner anderen Prüfungsbehörde vorgelegt und auch nicht veröffentlicht. \\[6ex]

\ort, den \today\\[3ex]


\rule[-0.2cm]{5cm}{0.5pt}

\textsc{\autor} 


%  Abstract ------------------------------------------------------------------------------------------
% !TEX encoding = UTF-8 Unicode
% !TEX root =  ../Arbeit.tex

\chapter*{Abstract}
\label{cha:Abstract}

\thispagestyle{empty}


Abstract

% Liste mit Todos -------------------------------------------------------------------------------------
\phantomsection
\markboth{Zu erledigende Punkte}{Liste der noch zu erledigenden Punkte}
\listoftodos
\newpage



% Seitennummerierung -----------------------------------------------------------------------------------
%   Vor dem Hauptteil werden die Seiten in großen römischen Ziffern 
%   nummeriert.
% ------------------------------------------------------------------------------------------------------
\ofoot{\pagemark \\[4ex]}
\pagenumbering{Roman}

% Inhaltsverzeichnis --------------------------------------------------------------------------------------
\phantomsection{} % Sorgt für korrekte Aufnahme des Inhaltsverzeichnisses in das Inhaltsverzeichnis
\addcontentsline{toc}{chapter}{Inhaltsverzeichnis}
\tableofcontents{}


%%
% Anzeigen des Abkürzungsverzeichnisses und des Glossars -------------------------------------
%%

% Abkürzungsverzeichnis anzeigen
\newpage
\markboth{Abkürzungsverzeichnis}{Abkürzungsverzeichnis}
\printglossary[style=listdotted, type=\acronymtype, title={Abkürzungsverzeichnis}, toctitle={Abkürzungsverzeichnis}]
\label{cha:Abkuerzungsverzeichnis}

% print glossary
\newpage
\printglossary[style=altlist,type=main]
\label{cha:Glossar}


% Restliche Verzeichnisse ------------------------------------------------------------------------------
\listoffigures{} % Abbildungsverzeichnis
\listoftables{} % Tabellenverzeichnis
\renewcommand{\lstlistlistingname}{Verzeichnis der Listings}
\lstlistoflistings{} % Listings-Verzeichnis

% arabische Seitenzahlen im Hauptteil ------------------------------------------------------------------
\clearpage{}
\pagenumbering{arabic}


% die Inhaltskapitel werden in "Inhalt.tex" inkludiert -------------------------------------------------
% !TEX encoding = UTF-8 Unicode
% !TEX root =  Arbeit.tex

% Hier können die einzelnen Kapitel inkludiert werden. Sie müssen in den 
% entsprechenden .TEX-Dateien vorliegen. Die Dateinamen können natürlich 
% angepasst werden.

% !TEX encoding = UTF-8 Unicode
% !TEX root =  ../Paper.tex

\chapter{Einleitung}
\label{cha:Einleitung}


\section{Lorem Ipsum}
\label{sec:ZielDerArbeit}

Lorem ipsum dolor sit amet, consetetur \gls{Glossareintrag} \ac{AGPL} sadipscing elitr, sed diam nonumy eirmod tempor invidunt ut labore et dolore magna aliquyam erat, sed diam voluptua. At vero eos et accusam et justo duo dolores et ea rebum. Stet clita kasd gubergren, no sea takimata sanctus est Lorem ipsum dolor sit amet. Lorem ipsum dolor sit amet, consetetur sadipscing elitr, sed diam nonumy eirmod \ac{AGPL} tempor invidunt ut labore et dolore magna aliquyam erat, sed diam voluptua. At vero eos et accusam et justo duo dolores et ea rebum. Stet clita kasd gubergren, no sea takimata sanctus est Lorem ipsum dolor sit amet. Lorem ipsum dolor sit amet, consetetur sadipscing elitr, sed diam nonumy eirmod tempor invidunt ut labore et dolore magna aliquyam erat, sed diam voluptua. At vero eos et accusam et justo duo dolores et ea rebum. Stet clita kasd gubergren, no sea takimata sanctus est Lorem ipsum dolor sit amet. \citep{Testautor:2002fk}

\section{Duis Autem}
\label{sec:Duis}

Duis autem vel eum iriure dolor in hendrerit in vulputate velit esse molestie consequat, vel illum dolore eu feugiat nulla facilisis at vero eros et accumsan et iusto odio dignissim qui blandit praesent luptatum zzril delenit augue duis dolore te feugait nulla facilisi. Lorem ipsum dolor sit amet, consectetuer adipiscing elit, sed diam nonummy nibh euismod tincidunt ut laoreet dolore magna aliquam erat volutpat. 

\section{Ut Wisi}
\label{sec:Ut}

Ut wisi enim ad minim veniam, quis nostrud exerci tation ullamcorper suscipit lobortis nisl ut aliquip ex ea commodo consequat. Duis autem vel eum iriure dolor in hendrerit in vulputate velit esse molestie consequat, vel illum dolore eu feugiat nulla facilisis at vero eros et accumsan et iusto odio dignissim qui blandit praesent luptatum zzril delenit augue duis dolore te feugait nulla facilisi. 



\clearpage{}
\pagenumbering{Roman}
\setcounter{page}{7} %%% Dieser Pagecounter muss entsprechend der verbrauchten Seiten im Inhaltsverzeichnis angepasst werden. Endet das IHV bei Seite III, so muss hier 4 eingetragen werden


% Literaturverzeichnis anzeigen
\phantomsection
\addcontentsline{toc}{chapter}{\bibliographyheading}

\setlength\bibitemsep{2\itemsep} % Abstand zwischen einträgen verdoppeln
\setlength{\bibhang}{0.2cm}

% Fette Literaturverzeichnis einträge
\xpretobibmacro{author}{\mkbibbold\bgroup}{}{}
\xapptobibmacro{author}{\egroup}{}{}
\xpretobibmacro{bbx:editor}{\mkbibbold\bgroup}{}{}
\xapptobibmacro{bbx:editor}{\egroup}{}{}

\renewcommand*{\labelnamepunct}{\mkbibbold{\addcolon\space}} 

\markboth{\bibliographyheading}{\bibliographyheading}

\printbibliography[heading=literaturverzeichnis]


% Index ------------------------------------------------------------------------------------------------
%   Zum Erstellen eines Index, die folgende Zeile auskommentieren.
% ------------------------------------------------------------------------------------------------------
%\printindex


% Anhang -----------------------------------------------------------------------------------------------
%   Die Inhalte des Anhangs werden analog zu den Kapiteln inkludiert.
%   Dies geschieht in der Datei "Anhang.tex".
% ------------------------------------------------------------------------------------------------------
\begin{appendix}
    \clearpage{}
    \pagenumbering{roman}
    \chapter{Anhang}
    \label{sec:Anhang}
    % Rand der Aufzählungen in Tabellen anpassen
    \setdefaultleftmargin{1em}{}{}{}{}{}
    % !TEX encoding = UTF-8 Unicode
% !TEX root =  Arbeit.tex

\section{Beiliegende CD}
\label{sec:BeiliegendeCd}

\subsection{Inhaltsverzeichnis der CD}
\label{subsec:InhaltsverzeichnisCd}

\begin{enumerate}
	\item Die gesamte Bachelorarbeit als PDF-Datei
	\item Alle verwendeten Online-Quellen als PDF-Ausdruck
	\item Sonstige Quelltexte
	\item Alle in der Bachelorarbeit verarbeiteten Diagramme
	\item Alle in der Bachelorarbeit verarbeiteten Grafiken
	\item Alle in der Bachelorarbeit verarbeiteten Screenshots
\end{enumerate}
\end{appendix}

\end{document}
